\chapter{Обзор литературы}

\section{Методы интерполяции видеокадров}

Интерполяция видеокадров (Video Frame Interpolation, VFI) представляет собой задачу синтеза промежуточных кадров между двумя исходными для повышения частоты кадров и плавности видеопоследовательности~\cite{kye2026acevfi}. Данная технология находит широкое применение в различных областях: от создания эффекта замедленного движения и повышения качества видеоконтента до сжатия видео и синтеза новых ракурсов~\cite{huang2022}.

С развитием глубокого обучения методы VFI эволюционировали от традиционных подходов, основанных на компенсации движения, к сложным нейросетевым архитектурам~\cite{kye2026acevfi}. 
Современные методы интерполяции можно классифицировать на несколько категорий в зависимости от архитектурных принципов~\cite{kye2026acevfi}.

\subsection*{Методы на основе оптического потока (Flow-based Methods)}

Методы данной категории основываются на явной оценке оптического потока~---~плотного поля векторов, описывающего смещение каждого пикселя между кадрами~\cite{IFRNet2022}.

Графически оптический поток представляется в виде векторного поля, где каждому пикселю сопоставлен вектор смещения. 
Пример такой визуализации показан на~рис.~\ref{img:flow}.
\includeimage
	{flow}
	{f}
	{H}
	{\textwidth}
	{Визуализация оптического потока $\mathbf{V}_{0\to1}$ как векторного поля.
		Стрелки на рис. (c), наложенные на кадр $I_0$ из (a), показывают смещение каждого пикселя к его положению в кадре $I_1$ (рис. (b)). 
		Длина стрелки пропорциональна величине смещения, направление~---~траектории движения.}
	
Типичный конвейер таких методов включает три основных этапа:

\begin{enumerate}
	\item \textbf{оценка потока}: расчет оптического потока между входными кадрами $I_0$ и $I_1$ с использованием нейросетевых архитектур, таких как FlowNet~\cite{ilg2017}, PWC-Net~\cite{sun2018} или FastFlowNet. Современные подходы часто применяют пирамидальные представления признаков для обработки больших перемещений~\cite{sun2018};
	\item \textbf{деформация (warping)}: использование рассчитанных полей потока для пространственного выравнивания пикселей или признаков входных кадров относительно позиции целевого промежуточного кадра, как правило, с использованием обратного деформирования (backward warping)~\cite{jiang2018};
	\item \textbf{синтез и уточнение}: поскольку деформированные кадры содержат артефакты (разрывы, окклюзии), применяется отдельная сеть синтеза для восстановления целевого кадра~\cite{kye2026acevfi}.
\end{enumerate}

Ключевым преимуществом потоковых методов является явное моделирование движения, обеспечивающее интерпретируемость и возможность интерполяции в произвольные моменты времени~\cite{huang2022}. 
Однако качество таких подходов критически зависит от точности оценки потока и значительно страдает при окклюзиях, больших перемещениях и изменении освещения~\cite{IFRNet2022}.

\subsection*{Методы на основе адаптивных ядер (Kernel-based Methods)}

В отличие от потоковых подходов, методы на основе адаптивных ядер не оценивают движение явно~\cite{niklaus2017}.
Вместо этого они используют архитектуру энкодер-декодер для извлечения признаков из входных кадров и генерации весов для локальных сверточных ядер, которые непосредственно применяются для деформации входных кадров~\cite{niklaus2017}. 
Классическими представителями данного подхода являются SepConvNet~\cite{niklaus2017}, AdaCoFNet и EDSC.
Основное преимущество таких методов~---~устойчивость в слаботекстурных областях, однако они сложны в настройке для произвольной временной интерполяции и диагностики ошибок~\cite{kye2026acevfi}.

\subsection*{Гибридные методы (Hybrid Methods)}

Гибридные подходы комбинируют преимущества потоковых и kernel-методов, используя оценку потока для глобального выравнивания и адаптивные ядра для уточнения локальных деталей~\cite{kye2026acevfi}. 
Это позволяет достигать более высокого качества, но зачастую сопровождается увеличением вычислительной сложности~\cite{huang2022}. 
Примером такого подхода является модель DAIN, использующая глубину для улучшения интерполяции.

\subsection*{Методы на основе трансформеров (Transformer-based Methods)}

В последние годы значительное внимание привлекают методы на основе трансформеров, которые используют механизмы глобального внимания для захвата долгосрочных пространственно-временных зависимостей между пикселями кадров~\cite{kye2026acevfi}. 
В отличие от сверточных сетей, трансформеры способны эффективно моделировать большие перемещения, поскольку операция внимания не ограничена локальным рецептивным полем. 
Среди таких подходов можно выделить VideoINR, использующий трансформеры для неявного представления видео, и TTVFI, комбинирующий трансформеры с модулями уточнения движения. 
Несмотря на высокое качество синтеза, вычислительная сложность трансформеров, квадратично зависящая от размера входного патча, остается существенным ограничением для их применения в реальном времени~\cite{kye2026acevfi}.

\section{Архитектура IFRNet}

Особый интерес для данной работы представляет архитектура \texttt{IFRNet} (Intermediate Feature Refine Network)~\cite{IFRNet2022}. 
Её ключевая инновация заключается в объединении этапов оценки потока и уточнения признаков в единую структуру энкодер-декодер, что позволяет достичь взаимного улучшения этих двух элементов и делает модель лёгкой и быстрой~\cite{IFRNet2022}.

\subsection*{Ключевая идея}

Существующие потоковые методы VFI, достигающие высокого качества, как правило, следуют конвейеру, разделяющему оценку промежуточного потока и синтез изображения на отдельные модули~\cite{IFRNet2022}. 
Это приводит к двум основным недостаткам~\cite{IFRNet2022}:
\begin{enumerate}
	\item разделение промежуточного потока и контекстных признаков на отдельные энкодер-декодеры игнорирует возможность их взаимного улучшения;
	\item каскадная архитектура существенно увеличивает задержку вывода и количество параметров модели, ограничивая применение в реальном времени и на мобильных устройствах~\cite{IFRNet2022}.
\end{enumerate}

\texttt{IFRNet} впервые объединяет эти два этапа в единую модель энкодер-декодер, где промежуточный поток и промежуточные признаки совместно уточняются, помогая друг другу~\cite{IFRNet2022}.
На рис.~\ref{img:diff_arch} представлены различные подходы VFI, основанные на потоке.
\includeimage
	{diff_arch}
	{f}
	{H}
	{\textwidth}
	{Различные подходы VFI, основанные на потоке~\cite{IFRNet2022}}

\subsection*{Структура сети}

\subsubsection*{Пирамидальный энкодер}

Сеть принимает на вход два кадра $I_0$ и $I_1$. 
Общий (shared) энкодер извлекает пирамиду признаков из каждого кадра. 
Энкодер построен из блоков двух сверток $3\times3$ на каждом уровне пирамиды (с шагами 2 и 1 соответственно). 
\texttt{IFRNet} извлекает 4 уровня пирамиды признаков, что составляет 8 сверточных слоев, каждый из которых сопровождается активацией PReLU. 
Количество каналов увеличивается с уровнем: 32, 48, 72 и 96, формируя признаки $\phi_0^k, \phi_1^k$ для каждого уровня $k \in \{1,2,3,4\}$~\cite{IFRNet2022}.

\subsubsection*{Декодеры <<от грубого к точному>>}

Декодеры последовательно обрабатывают признаки от самого грубого (низкое разрешение) уровня к самому детальному. 
Каждый декодер $\mathcal{D}^k$ одновременно предсказывает два компонента~\cite{IFRNet2022}:

\begin{itemize}
	\item \textbf{двусторонний промежуточный поток} ($F_{t \to 0}^k, F_{t \to 1}^k$): векторы движения от целевого момента времени к первому и второму входным кадрам;
	\item \textbf{реконструированный промежуточный признак} ($\hat{\phi}_t^k$): представление целевого кадра в пространстве признаков.
\end{itemize}

Схема архитектуры и декодера представлены на рисунках~\ref{img:ifrnet_arch.png}~-~\ref{img:ifrnet_arch.pdf}.

\includeimage
{ifrnet_arch.png}
{f}
{H}
{.8\textwidth}
{Обзор архитектуры и функций потери данных в \texttt{IFRNet}~\cite{IFRNet2022}}

\includeimage
{ifrnet_arch.pdf}
{f}
{H}
{.8\textwidth}
{Подробная схема декодера на каждом уровне пирамиды~\cite{IFRNet2022}}

Процесс устроен циклически: более точный поток позволяет лучше деформировать исходные пирамидальные признаки посредством обратного деформирования (backward warping), порождая более качественные $\tilde{\phi}_0^k, \tilde{\phi}_1^k$, что дает лучший материал для реконструкции $\hat{\phi}_t^k$. 
В свою очередь, наличие качественного промежуточного признака предоставляет недостающую опорную информацию для более точного предсказания движения на следующем уровне~\cite{IFRNet2022}.

Каждый декодер представляет собой компактную сеть из шести сверток $3\times3$ и одной деконволюции $4\times4$ (с шагом $1/2$). 
Для сохранения эффективности при достаточно большом рецептивном поле третья и пятая свертки обновляют только часть каналов. 
Остаточные соединения способствуют распространению информации и совместному уточнению~\cite{IFRNet2022}.

Благодаря постепенно уточняемому промежуточному признаку, содержащему информацию об окклюзиях и глобальном контексте, \texttt{IFRNet} может генерировать маску слияния и компенсировать детали движения без необходимости в дополнительном модуле синтеза или уточнения, в отличие от многих других методов~\cite{IFRNet2022}.

\subsection*{Синтез кадра}

На финальном этапе выходные карты признаков преобразуются в изображение. 
Сеть генерирует маску слияния $M$ и остаток $R$. Итоговый кадр $\hat{I}_t$ синтезируется по формуле:

\begin{align*}
	\hat{I}_t &= M \odot \tilde{I}_0 + (1 - M) \odot \tilde{I}_1 + R, \\
	\tilde{I}_0 &= w(I_0, F_{t\to0}), \quad \tilde{I}_1 = w(I_1, F_{t\to1}),
\end{align*}

где $\tilde{I}_0$ и $\tilde{I}_1$~---~деформированные входные кадры. Маска $M$ определяет вклад каждого кадра (решая проблему окклюзий), а остаток $R$ добавляет детали текстуры, утраченные при деформации. 
Маска слияния адаптивно усредняет области, видимые в обоих кадрах, а остаток обычно активируется на границах движения и контурах объектов~\cite{IFRNet2022}.

\subsection*{Функции потерь}

Для обучения \texttt{IFRNet} используются специальные функции потерь, позволяющие эффективно обучать модель~\cite{IFRNet2022}:

\begin{itemize}
	\item \textbf{задаче-ориентированный дистилляционный лосс} (task-oriented optical flow distillation loss): этот подход адаптивно регулирует в пространстве робастность дистилляционного знания от более тяжелой предобученной teacher-сети, фокусируясь на изучении полезной информации для синтеза кадра и игнорируя потенциально вредный шум~\cite{IFRNet2022};
	\item \textbf{геометрический консистентный лосс} (feature space geometry consistency loss): этот лосс использует признаки из эталонного промежуточного кадра для регуляризации реконструированных промежуточных признаков, сохраняя правильную геометрическую структуру объектов~\cite{IFRNet2022}.
\end{itemize}

Такой подход позволяет \texttt{IFRNet} достигать качества на уровне SOTA по метрикам \texttt{PSNR}/\texttt{SSIM} при высокой скорости работы и малом числе параметров, так как вся логика упакована в один эффективный блок~\cite{IFRNet2022}.

\section{Методы оптимизации нейросетевых архитектур}

Существует несколько подходов к снижению вычислительной сложности моделей глубокого обучения.

\textbf{Сокращение ширины сети}~---~уменьшение числа каналов в сверточных слоях позволяет значительно сократить число параметров и FLOPs, так как вычислительная стоимость сверточного слоя с ядром $k \times k$ приближённо пропорциональна $H W k^2 C_{in} C_{out}$. 
Одновременное уменьшение входных и выходных каналов в два раза теоретически снижает стоимость примерно в четыре раза.

\textbf{Использование динамических сверток}~---~данный подход позволяет адаптивно применять вычислительные ресурсы только к значимым пикселям изображения. 
В работе CALC-VFS~\cite{giuliani2023calc} предлагается внедрение динамических сверток в архитектуру \texttt{IFRNet}, что позволяет достичь снижения вычислительной сложности на 20-40\% при сохранении схожего качества.

\textbf{Переход к альтернативным цветовым пространствам}~---~использование форматов с субдискретизацией цветоразностных компонент (например, YUV420) позволяет сократить объем входных данных в два раза по сравнению с RGB, что потенциально снижает вычислительную нагрузку на этапе обработки цветовых каналов.


\subsection*{Уменьшение числа признаков в два раза}

Данный метод заключается в замене ширины каналов $C_l$ на

\begin{equation}
	C_l' = \left\lfloor \frac{C_l}{2} \right\rfloor
\end{equation}

для каждого уровня пирамиды. Такая модификация затрагивает энкодер, refine-блоки и часть синтеза. 
Архитектурная логика IFRNet при этом сохраняется: остаются те же уровни пирамиды, тот же порядок уточнения от грубого масштаба к детальному, те же потоки, маски и residual-компоненты.
Архитектурная логика \texttt{IFRNet} при этом сохраняется: остаются те же уровни пирамиды, тот же порядок уточнения от грубого масштаба к детальному, те же потоки, маски и residual-компоненты.

Главное преимущество метода~---~простота внедрения. 
Модель по-прежнему принимает два кадра и предсказывает промежуточный кадр, но использует более узкое внутреннее представление. 
Ожидаемый эффект~---~уменьшение числа параметров, FLOPs и объёма промежуточных активаций.

Недостаток метода связан с уменьшением представительной способности сети. Более узкие признаки могут хуже кодировать мелкие текстуры, сложные движения и области с частичными окклюзиями.

\subsection*{Переход к YUV420}

В цветовом пространстве RGB все три канала имеют одинаковое пространственное разрешение, хотя человеческое зрение более чувствительно к яркости, чем к цветоразностной информации. 
Это свойство широко используется в системах видеокомпрессии: стандарты сжатия, такие как HEVC и H.264, ожидают входные данные в форматах YUV с субдискретизацией цветности.

В YUV-представлении изображение разделяется на яркостную компоненту $Y$ и две цветоразностные компоненты $U$ и $V$. 
Формат YUV420 хранит яркостную компоненту $Y$ в полном разрешении, а цветоразностные компоненты $U$ и $V$~---~с понижением разрешения в два раза по горизонтали и вертикали. 
Для кадра размера $H \times W$ RGB содержит $3HW$ значений, тогда как YUV420:
\begin{equation}
HW + \frac{HW}{4} + \frac{HW}{4} = 1.5HW
\end{equation}

значений, что в два раза меньше данных по сравнению с RGB. Это делает YUV420 перспективным форматом для ускорения интерполяции, поскольку объём входных данных существенно сокращается без потери информации, критичной для зрительного восприятия.

\includeimage
{yuv}
{f}
{H}
{.9\textwidth}
{Фотография и её YUV-компоненты}

\subsubsection*{Стратегии применения YUV420 в IFRNet}

В контексте архитектуры IFRNet возможны две стратегии использования YUV420:

\begin{enumerate}
	\item \textbf{раздельная интерполяция компонент}: интерполяция компонент $Y$, $U$ и $V$ производится раздельно, причём для цветоразностных компонент $U$ и $V$ используется сеть с уменьшенным разрешением или отдельные ветви. 
	Такой подход позволяет сократить вычисления на цветовых каналах, так как они обрабатываются на четверти исходного разрешения.
	
	\item \textbf{интерполяция на основе яркости}: оценка движения выполняется по яркостной компоненте $Y$, а цветоразностные компоненты переносятся с использованием найденного поля движения. 
	Этот подход основан на наблюдении, что движение объектов одинаково для яркости и цвета, а цветовая информация более устойчива к небольшим пространственным ошибкам.
\end{enumerate}

Второй подход является более эффективным, поскольку позволяет избежать дополнительных вычислений на цветовых компонентах.

\subsubsection*{Преимущества и риски}

\textbf{Преимущества} перехода к YUV420:
\begin{itemize}
	\item сокращение объёма входных данных в два раза по сравнению с RGB;
	\item потенциальное снижение вычислительной нагрузки на этапе обработки цветовых каналов;
	\item соответствие стандартам видеосжатия, что упрощает интеграцию в существующие системы;
	\item возможность использования аппаратных ускорителей, оптимизированных для работы с YUV420.
\end{itemize}

\textbf{Основные риски} связаны с возможным появлением цветовых артефактов:
\begin{itemize}
	\item неточное восстановление цветоразностных компонент может приводить к смещению оттенков;
	\item размытие цветных границ из-за пониженного разрешения $U$ и $V$;
	\item ошибки в насыщенных областях и на резких цветовых переходах;
	\item потенциальная потеря мелких цветных деталей.
\end{itemize}

\section{Метрики оценки качества}

Оценка методов интерполяции видеокадров проводится с помощью комбинации метрик, измеряющих различные аспекты качества~\cite{kye2026acevfi}. 
Выбор метрик определяется необходимостью комплексного анализа: пиксельные метрики позволяют оценить численную точность восстановления, тогда как перцептуальные метрики отражают визуальное качество, согласованное с человеческим восприятием~\cite{kye2026acevfi}. 
В данном исследовании используются четыре метрики: \texttt{PSNR}, \texttt{SSIM}, \texttt{LPIPS} и \texttt{DISTS}, что позволяет всесторонне сравнить эффективность предлагаемых методов оптимизации.

\subsection*{Пиксельные метрики}

Пиксельные метрики сравнивают интерполированный кадр $\hat{I}_t$ с эталонным $I_t$ на уровне отдельных пикселей или их локальных статистик.

\subsubsection*{PSNR (Peak Signal-to-Noise Ratio)}

\texttt{PSNR} является наиболее распространённой метрикой оценки качества восстановления изображений. Она вычисляется на основе среднеквадратичной ошибки (Mean Squared Error, MSE) между пикселями эталонного и восстановленного изображений:

\begin{equation}
	\text{MSE} = \frac{1}{HW} \sum_{i=1}^{H} \sum_{j=1}^{W} \left( I_t(i,j) - \hat{I}_t(i,j) \right)^2,
\end{equation}

\begin{equation}
	\text{PSNR} = 10 \cdot \log_{10}\left(\frac{MAX_I^2}{\text{MSE}}\right),
\texttt{PSNR} = 10 \cdot \log_{10}\left(\frac{MAX_I^2}{\text{MSE}}\right),
\end{equation}

где $H$ и $W$~---~высота и ширина изображения, $MAX_I$~---~максимальное значение пикселя (для 8-битных изображений $MAX_I = 255$). Высокие значения \texttt{PSNR} указывают на высокое численное сходство между изображениями.

Преимуществом \texttt{PSNR} является её простота, физический смысл и возможность легко сравнивать результаты с другими работами благодаря широкой распространённости. Однако главный недостаток заключается в том, что \texttt{PSNR} плохо коррелирует с человеческим восприятием: два изображения могут иметь одинаковый \texttt{PSNR}, но одно будет выглядеть отлично, а другое~---~содержать заметные артефакты, поскольку \texttt{PSNR} не учитывает структурную информацию и пространственную организацию пикселей.

\subsubsection*{SSIM (Structural Similarity Index)}

\texttt{SSIM} была разработана для преодоления недостатков \texttt{PSNR} и оценки сходства изображений с учётом особенностей зрительной системы человека. Метрика вычисляет сходство по трём компонентам:

\begin{equation}
	\text{SSIM}(x, y) = \frac{(2\mu_x\mu_y + C_1)(2\sigma_{xy} + C_2)}{(\mu_x^2 + \mu_y^2 + C_1)(\sigma_x^2 + \sigma_y^2 + C_2)},
\texttt{SSIM}(x, y) = \frac{(2\mu_x\mu_y + C_1)(2\sigma_{xy} + C_2)}{(\mu_x^2 + \mu_y^2 + C_1)(\sigma_x^2 + \sigma_y^2 + C_2)},
\text{SSIM}(x, y) = \frac{(2\mu_x\mu_y + C_1)(2\sigma_{xy} + C_2)}{(\mu_x^2 + \mu_y^2 + C_1)(\sigma_x^2 + \sigma_y^2 + C_2)},
\end{equation}

где $x$ и $y$~---~соответствующие локальные блоки эталонного и восстановленного изображений; $\mu_x$, $\mu_y$~---~средние значения; $\sigma_x$, $\sigma_y$~---~дисперсии; $\sigma_{xy}$~---~ковариация; $C_1$ и $C_2$~---~константы, вводимые для стабилизации деления. Итоговое значение \texttt{SSIM} вычисляется как среднее по всем локальным блокам изображения.

\texttt{SSIM} значительно лучше согласуется с человеческим восприятием, чем \texttt{PSNR}, и стала стандартом де-факто в области обработки изображений. Однако метрика имеет ограничения: она может давать нестабильные результаты на изображениях с низкой контрастностью или яркостью, а также чувствительна к мелким пространственным сдвигам и деформациям.

\subsection*{Перцептуальные метрики}

Перцептуальные метрики оценивают визуальное качество через признаки, извлечённые из глубоких нейросетей, что позволяет приблизить оценку к субъективному восприятию человека~\cite{kye2026acevfi}.

\subsubsection*{LPIPS (Learned Perceptual Image Patch Similarity)}

\texttt{LPIPS} использует предобученную свёрточную нейросеть (как правило, VGG) для извлечения карт признаков из промежуточных слоёв. Метрика вычисляет взвешенное расстояние между признаками эталонного и восстановленного изображений на каждом уровне:

\begin{equation}
	\text{LPIPS}(I_t, \hat{I}_t) = \sum_{l} \sum_{c} \lambda_{l,c} \cdot \| \phi_{l,c}(I_t) - \phi_{l,c}(\hat{I}_t) \|_2,
\end{equation}

где $\phi_{l,c}$~---~карта признаков на слое $l$ для канала $c$, а $\lambda_{l,c}$~---~изученные веса, оптимизированные для согласования с субъективными оценками человека. Низкие значения \texttt{LPIPS} указывают на лучшее перцептуальное качество.

\texttt{LPIPS} демонстрирует более высокую корреляцию с человеческими оценками по сравнению с \texttt{PSNR} и \texttt{SSIM}, что делает её предпочтительной для задач, где визуальное качество важнее численной точности. Однако метрика не лишена недостатков: она может быть обманута специальными атаками, и её работа не всегда интерпретируема. Кроме того, \texttt{LPIPS} чувствительна к глобальным изменениям цвета и контраста.

\subsubsection*{DISTS (Deep Image Structure and Texture Similarity)}

\texttt{DISTS} является более современной метрикой, разработанной для раздельной оценки структурного и текстурного сходства. В отличие от \texttt{LPIPS}, \texttt{DISTS} вычисляет две компоненты:

\begin{enumerate}
	\item \textbf{Структурное сходство} (Structure Similarity)~---~оценивает пространственное расположение признаков, то есть точную структуру объектов;
	\item \textbf{Текстурное сходство} (Texture Similarity)~---~оценивает статистики признаков по пространству, что позволяет быть устойчивым к передискретизации текстур.
\end{enumerate}

Математически метрика может быть представлена как:

\begin{equation}
	\text{DISTS}(I_t, \hat{I}_t) = \sum_{l} \sum_{c} \lambda_{l,c} \cdot \left\| \text{corr}\left(\phi_{l,c}(I_t)\right) - \text{corr}\left(\phi_{l,c}(\hat{I}_t)\right) \right\|_2,
\end{equation}

где $\text{corr}$ обозначает вычисление корреляционной структуры карты признаков, а $\lambda_{l,c}$~---~изученные веса.

Ключевая особенность \texttt{DISTS} заключается в том, что она позволяет учесть вариативность текстур (например, при синтезе травы, воды или песка), где небольшие перестановки текстурных элементов не должны существенно влиять на оценку качества. Исследования показывают, что \texttt{DISTS} демонстрирует наивысшую корреляцию с субъективными оценками человека среди всех современных метрик и устойчива к небольшим геометрическим искажениям.

\subsection*{Обоснование выбора метрик}

Комплексное использование пиксельных и перцептуальных метрик позволяет объективно оценивать как численную точность, так и визуальное качество интерполированных кадров, что критически важно для сравнения методов оптимизации~\cite{kye2026acevfi}.

В контексте данной работы выбор четырёх метрик (\texttt{PSNR}, \texttt{SSIM}, \texttt{LPIPS}, \texttt{DISTS}) обусловлен следующими соображениями:
\t\item \textbf{\texttt{PSNR}} и \textbf{\texttt{SSIM}} позволяют оценить, насколько сильно изменились пиксели и простая структура изображения после применения методов оптимизации (сокращения каналов или перехода в YUV420);

\begin{itemize}
	\item \textbf{\texttt{PSNR}} и \textbf{\texttt{SSIM}} позволяют оценить, насколько сильно изменились пиксели и простая структура изображения после применения методов оптимизации (сокращения каналов или перехода в YUV420);
	\item \textbf{\texttt{LPIPS}} и \textbf{\texttt{DISTS}} показывают, насколько эти изменения заметны для человеческого глаза, что критически важно, поскольку цель оптимизации~---~сохранить визуальное качество при повышении вычислительной эффективности.
\end{itemize}

Например, если после оптимизации значение \texttt{PSNR} снизилось, но \texttt{LPIPS} и \texttt{DISTS} изменились незначительно, это означает, что выигрыш в скорости достигнут без заметной для наблюдателя потери качества.